% PhyriadFG — por qué la captura externa: lo que recibe un FG dentro del motor vs lo que recibe PhyriadFG
% Build (MiKTeX, figure contract v1, portfolio/figure-contract/README.md §1):
%   latex -interaction=nonstopmode phyriadfg_por_que_captura_externa.tex
%   dvisvgm --no-fonts --bbox=papersize -o phyriadfg_por_que_captura_externa.svg phyriadfg_por_que_captura_externa.dvi 2> dvisvgm.log
%   (any "failed to create" in dvisvgm.log = FAILED build)
%   svgo --multipass phyriadfg_por_que_captura_externa.svg -o phyriadfg_por_que_captura_externa.min.svg
% Colours: phyriad-figure-colors.sty (copy it next to this file). Ink = TeX black; ONE accent (figAccent) on
% what PhyriadFG estimates itself (the flow box of thread F); figMuted for the thread/GPU notes.
% Redrawn 2026-10-10 with the dossier's Anexo corrections 1, 2, 3 and 4 (PORTFOLIO_DOSSIER.md).
% Sources:
%   top row    - AMD GPUOpen, "FidelityFX Frame Interpolation 1.1.3",
%                https://gpuopen.com/manuals/fidelityfx_sdk/techniques/frame-interpolation/, re-read 2026-10-10 via
%                a summarising fetch [V2]: inputs = presentation colour buffer(s), dilated depth, dilated motion
%                vectors, camera near/far/FOV, frame-time delta, AND "the optical flow motion vectors"; it is
%                "designed to work together with the FfxOpticalFlow"; algorithm stage "Compute optical flow
%                vector field".
%   bottom row - public README.md:3-5 ("with no access to the game engine, its motion vectors, or its depth"),
%                README.md:99-100 (disocclusion is the hard case), docs/ARCHITECTURE.md:22 (thread C acquires
%                the frame by DDA or WGC; iGPU converts+packs), :23 (thread F: MV + SAD via a pyramid plus the
%                global motion model gme), :24 (thread P: warp to t, blend, quality, present), :33 (pipeline).
% Made with my soul - Swately <3
\documentclass[border=12pt]{standalone}
\usepackage[utf8]{inputenc}
\usepackage[T1]{fontenc}
\usepackage{lmodern}
\renewcommand{\familydefault}{\sfdefault}
\usepackage{amsmath}
\usepackage{phyriad-figure-colors}
\usepackage{tikz}
\usetikzlibrary{arrows.meta,positioning,calc}

\begin{document}
\begin{tikzpicture}[
  >=Latex, line width=0.6pt,
  box/.style={draw, rounded corners, minimum height=1.25cm, align=center, font=\small},
  data/.style={draw, dashed, rounded corners, minimum height=1.25cm, align=left, font=\footnotesize, inner sep=4pt},
  accentbox/.style={box, draw=figAccent, fill=figFillAccent, line width=1.1pt},
  row/.style={font=\small\bfseries, anchor=west},
  note/.style={font=\footnotesize, align=center, text=figMuted},
]

% ── fila inferior: PhyriadFG (se dibuja primero; la fila superior se alinea con ella) ──
\begin{scope}[yshift=-5.6cm]
\node[row] at (0,2.0) {PhyriadFG (captura externa)};
\node[box, text width=2.3cm] (game) at (1.3,0) {Juego\\(caja cerrada)};
\node[data, right=9mm of game, text width=4.4cm] (pin)
  {solo los píxeles finales\\de la pantalla\\(captura DDA o WGC)};
\node[accentbox, right=9mm of pin, text width=3.3cm] (flow)
  {Flujo óptico estimado\\de los píxeles (MV\,+\,SAD)\\+ modelo global};
\node[box, right=9mm of flow, text width=3.4cm] (warp)
  {Warp hacia atrás\\+ mezcla + capas de\\calidad (oclusión)};
\node[box, right=9mm of warp, text width=1.4cm] (pout) {$I_t$};
\draw[->] (game) -- (pin);
\draw[->] (pin) -- (flow);
\draw[->] (flow) -- (warp) node[midway, above, font=\footnotesize]{MV};
\draw[->] (warp) -- (pout);

% hilos y GPU bajo cada etapa
\node[note, anchor=north] at (pin.south)  {hilo C: captura;\\la GPU integrada convierte y empaqueta};
\node[note, anchor=north] at (flow.south) {hilo F\\segunda GPU};
\node[note, anchor=north] at (warp.south) {hilo P\\GPU principal};
\end{scope}

% ── fila superior: FG dentro del motor ──────────────────────────────────
\node[row] at (0,2.15) {FG dentro del motor (p.\,ej. FSR 3 Frame Interpolation)};
\node[box, text width=2.3cm] (eng) at (game.center |- 0,0) {Motor del juego};
\node[data, text width=4.4cm] (ein) at (pin.center |- 0,0)
  {color del cuadro\\profundidad (depth)\\vectores de movimiento del motor\\cámara: near/far, FOV, $\Delta t$};
\node[box, text width=3.3cm] (eflow) at (flow.center |- 0,0) {Flujo óptico propio\\(\texttt{FfxOpticalFlow})};
\node[box, text width=3.4cm] (eint) at (warp.center |- 0,0) {Interpolador\\(dentro del motor)};
\node[box, text width=1.4cm] (eout) at (pout.center |- 0,0) {$I_t$};
\draw[->] (eng) -- (ein);
\draw[->] (ein) -- (eflow);
\draw[->] (eflow) -- (eint) node[midway, above, font=\footnotesize]{flujo};
\draw[->] (eint) -- (eout);
% profundidad, vectores y cámara también llegan directamente al interpolador
\draw[->] (ein.north) |- ($(eflow.north)+(0,0.55)$) node[pos=0.75, above, font=\footnotesize]{profundidad, vectores y cámara} -| (eint.north);

% ── lo que falta: profundidad y vectores del motor ─────────────────────
\draw[dashed] (ein.south) -- node[pos=0.45, font=\large, inner sep=1pt] {$\times$}
  node[pos=0.45, right, font=\footnotesize, align=left, xshift=4mm]
  {no llegan profundidad ni vectores del motor:\\PhyriadFG estima todo el movimiento solo de la imagen}
  (pin.north);

\end{tikzpicture}
\end{document}
